\SourceLeaf{380}%
lang gestört, und eine vorher \textit{nie} vorgekommene bedeutende Wandelbarkeit des Collimationsfehlers bemerklich. Die Quelle davon fand sich, nachdem sie vorher vergeblich in mancherlei andern Umständen gesucht war, in einem Knötchen des Fadens, welcher um die den Verschluss der Libelle sichernde Blasenhaut gebunden war, und, ein klein wenig zu dick, die innere Fläche der Hülse berührte: nachdem dieses Knötchen weggeschnitten war, so dass die Glasröhre bloss die nur wenig vortretenden Schraubenspitzen berührte, war die Beständigkeit des Collimationsfehlers, und die frühere schöne Harmonie aller Beobachtungen sogleich wieder hergestellt. Auch bei den in Frage stehenden Turiner Beobachtungen bemerken wir bedeutende Wandelbarkeit in dem Collimationsfehler, wir meinen nicht die grösseren Veränderungen von mehreren Minuten, die ohne Zweifel ihren guten dem Astronomen bekannten, obwohl bei den Beobachtungen nicht angeführten Grund gehabt haben, sondern die kleinern, welche zufällig scheinen. Gegenwärtig, wo man ein so vortreffliches Mittel hat, den Collimationsfehler jeden Augenblick ohne Umlegen zu bestimmen, wird die Auffindung der Ursache von ähnlichen Anomalien um so mehr erleichtert.

\begingroup\rightskip=0pt plus 1em
Im vierten Abschnitt wird der Anschluss des Mont Cenis an das Dreieckssystem vermittelst einer besondern Triangulirung und einer kleinen auf dem Plateau des Berges gemessenen Grundlinie, wie auch die astronomische Bestimmung des Azimuths der Verbindungslinie Mont Cenis --- Bellecombe mitgetheilt. Letztere ist zweimal gemacht; die Resultate der Jahre 1821, 1822, mit Repetitionskreisen von \textsc{Throughton} und \textsc{Reichenbach}, weichen \(8''6\) von einander ab, und man nahm, obgleich die spätere Bestimmung bei weitem zuverlässiger scheint, aus beiden das Mittel.\par
\endgroup

Eben so enthalten die beiden folgenden Abschnitte die astronomischen Bestimmungen der Azimuthe der Richtungslinien Mont Colombier --- Mont Granier und Turin neue Sternwarte --- Superga. In allen drei Fällen dienten die Meridianzeichen der resp. Mittagsfernröhre zur Grundlage dieser Bestimmungen. Endlich finden sich noch im sechsten Abschnitt die Operationen, durch welche eine früher von \textsc{Oriani} auf der Mailänder Sternwarte gemachte astronomische Azimuthalbestimmung auf die Orientirung der Seite in dem französischen Dreieckssystem Mailand Domthurm --- Busto übertragen wurde.

Im siebenten Abschnitte werden nun aus diesen ausgedehnten Operationen die Resultate für die Längengradmessung abgeleitet. Man bezog die Messungen
